\section{Introducción}
\label{sec:introduccion}

El filtrado adaptativo ocupa un lugar distinto al del filtrado clásico
(FIR/IIR de coeficientes fijos, abordado en trabajos previos de esta
misma serie) porque no exige conocer de antemano las características
estadísticas de la señal ni del sistema que se desea modelar o
compensar. En su lugar, el filtro ajusta sus propios coeficientes de
forma iterativa, muestra a muestra, a partir únicamente de la señal de
entrada y de una señal de error calculada en línea. Esta propiedad lo
hace especialmente adecuado para tres familias de problemas que
aparecen con frecuencia en instrumentación y comunicaciones: la
identificación ciega de sistemas desconocidos, la cancelación de ruido
y eco cuando no se dispone de un modelo explícito del canal de ruido,
y el seguimiento de sistemas cuyo comportamiento cambia en el tiempo.

El algoritmo \textbf{LMS} (\textit{Least Mean Squares}), introducido
por Widrow y Hoff a finales de la década de 1950, resuelve este
problema mediante una aproximación estocástica del descenso de
gradiente sobre el error cuadrático medio. En lugar de calcular el
gradiente exacto de $E[e^2[n]]$ —que requeriría conocer la estadística
conjunta de la señal, generalmente desconocida o variante—, LMS usa el
error instantáneo $e^2[n]$ como estimador insesgado de dicho gradiente,
obteniendo la regla de actualización:
\begin{equation}
w[n+1] = w[n] + \mu \, e[n] \, x[n],
\label{eq:lms}
\end{equation}
donde $x[n]$ es el vector de las $N$ muestras más recientes de la
entrada, $w[n]$ es el vector de pesos del filtro en el instante $n$,
$d[n]$ es la señal deseada, $y[n] = w[n]^\top x[n]$ es la salida del
filtro, y $e[n] = d[n] - y[n]$ es el error. El tamaño de paso $\mu$
controla cuánto se desplaza $w$ en cada iteración en la dirección que
reduce el error: un valor demasiado pequeño produce una convergencia
lenta pero segura, mientras que un valor demasiado grande acelera la
convergencia a costa de introducir oscilaciones crecientes que pueden
llevar el filtro a la inestabilidad numérica. Esta disyuntiva —el
llamado compromiso velocidad/estabilidad— es intrínseca a LMS y se
caracteriza formalmente mediante el mayor autovalor de la matriz de
autocorrelación de la entrada, como se detalla en la
Sección~\ref{sec:resultados}.

Una limitación práctica importante de LMS es que ese compromiso debe
resolverse eligiendo un único valor de $\mu$ antes de ejecutar el
filtro, bajo el supuesto implícito de que la potencia de la señal de
entrada es razonablemente estacionaria. Cuando la señal de entrada
cambia de potencia en el tiempo —como ocurre, por ejemplo, con la voz
humana, con señales de audio con pasajes fuertes y silencios, o con
canales de comunicación cuya ganancia varía—, un $\mu$ fijo deja de ser
óptimo en todo momento: si se eligió lo bastante pequeño para ser
estable en los segmentos de mayor potencia, resultará innecesariamente
lento en los segmentos de menor potencia, y viceversa.

El algoritmo \textbf{NLMS} (\textit{Normalized LMS}) resuelve
exactamente esta limitación normalizando el tamaño de paso por la
energía instantánea del vector de entrada:
\begin{equation}
w[n+1] = w[n] + \frac{\mu}{\varepsilon + \lVert x[n] \rVert^2} \, e[n] \, x[n],
\label{eq:nlms}
\end{equation}
donde $\varepsilon > 0$ es una constante pequeña que evita la división
por cero cuando la entrada es momentáneamente nula o muy débil. El
efecto de esta normalización es que el paso \emph{efectivo}
$\mu / (\varepsilon + \lVert x[n] \rVert^2)$ se autoajusta en cada
iteración a la potencia instantánea de la señal: cuando la entrada es
fuerte, el paso efectivo se reduce automáticamente para mantener la
estabilidad; cuando la entrada es débil, el paso efectivo crece para
no perder velocidad de convergencia. Esto permite usar, en la práctica,
un único parámetro $\mu$ (típicamente en el rango $0 < \mu \leq 2$,
con $\mu = 1$ como punto de partida habitual) que se comporta de forma
razonablemente estable e igualmente rápida sin importar cómo varíe la
potencia de la señal de entrada a lo largo del tiempo —una propiedad
que LMS, con su paso fijo, no puede ofrecer.

Este informe documenta dos implementaciones independientes, construidas
enteramente desde cero (sin recurrir a una función de filtrado
adaptativo ya empaquetada por una librería), que permiten verificar
experimentalmente estas propiedades teóricas en vez de darlas por
ciertas: \textbf{WORK\_ELEVEN} implementa y caracteriza LMS mediante
tres experimentos (identificación de sistema, cancelación adaptativa de
ruido y un barrido del tamaño de paso), y \textbf{WORK\_TWELVE}
implementa NLMS con los mismos tres experimentos más una comparación
directa, cabeza a cabeza, contra una reimplementación local de LMS bajo
una señal de entrada cuya potencia cambia abruptamente a mitad de la
prueba —el escenario exacto en el que la ventaja teórica de NLMS debería
manifestarse de forma medible.
